\begin{theorem}[Relative consistency of CH]
\label{thm:ch-relative-consistency}
If \(\ZF\) is consistent, then \(\ZFC+\mathsf{GCH}+(V=L)\) is consistent.
Hence if \(\ZFC\) is consistent, then \(\ZFC+\mathsf{CH}\) is consistent.
\end{theorem}

\begin{proof}
Suppose a derivation concludes \(\bottom\) from axioms
\(\sigma_1,\ldots,\sigma_k\) of \(\ZFC+\mathsf{GCH}+(V=L)\).  All the
\(\sigma_i\) and \(\bottom\) are sentences, so
\cref{lem:relativization-derivations} with \(m=0\) and \(\psi=\bottom\)
gives
\(\ZF\derives\sigma_1^{L}\land\cdots\land\sigma_k^{L}\to\bottom\), and the
relativization scheme cited in \cref{sec:constructibility-forcing-01}
gives \(\ZF\derives\sigma_i^{L}\) for each \(i\).  Conjunction
introduction and implication elimination yield \(\ZF\derives\bottom\),
contradicting the consistency of \(\ZF\).

For the second claim, every \(\ZF\)-derivation is a \(\ZFC\)-derivation,
so consistency of \(\ZFC\) gives consistency of \(\ZF\) and, by the first
claim, of \(\ZFC+\mathsf{GCH}+(V=L)\).  That theory derives
\(\mathsf{CH}\): universal elimination at the ordinal zero---formally,
existential elimination from the \(\ZFC\)-theorem that the ordinal zero
exists, followed by universal elimination and equality
substitutivity---produces the \(\alpha=0\) instance of the formalized
\(\mathsf{GCH}\) display, which is the formalized \(\mathsf{CH}\).  Now
suppose some derivation concludes \(\bottom\) with open assumptions in
\(\ZFC\cup\set{\mathsf{CH}}\).  At least one open occurrence of
\(\mathsf{CH}\) remains in it, since otherwise its open assumptions lie in
\(\ZFC\) and \(\ZFC\derives\bottom\), contrary to the assumed consistency
of \(\ZFC\).  Applying \(\neg I\) at the root, discharging every such
occurrence, gives \(\ZFC\derives\neg\mathsf{CH}\); monotonicity of the
derivability relation gives
\(\ZFC+\mathsf{GCH}+(V=L)\derives\neg\mathsf{CH}\); and negation
elimination against the derivation of \(\mathsf{CH}\) above gives
\(\ZFC+\mathsf{GCH}+(V=L)\derives\bottom\), contradicting the consistency
established in the first claim.  Hence \(\ZFC+\mathsf{CH}\) is
consistent.
\end{proof}

